\section{Keep RGBA values separate from display behavior}
Represent a decoded pixel as the numerical word \texttt{0xAARRGGBB}. The high byte is
alpha; the other three bytes retain the display convention from the earlier chapters.
In little-endian memory, the bytes of \texttt{0x80070605} are \texttt{05 06 07 80}.
In a QOI RGBA chunk, the component bytes instead appear as \texttt{07 06 05 80}.
The word's numerical fields, its memory bytes and the file's component order are three
different descriptions. Write down which one a particular inspection is showing.

The current MARS Bitmap Display builds an RGB image. It does not blend these words
against a background according to their high byte. A pixel with zero alpha can still
show its stored RGB values there. Keep alpha in the decoded words for the encoder;
do not infer its value from the display. This chapter does not implement compositing
or convert between linear and gamma-encoded component values.

When loading a complete color word into a 64-bit register, \texttt{lwu} provides the
zero-extended value. A signed \texttt{lw} can make the register's upper half all ones
when alpha sets bit thirty-one. That sign extension does not change a subsequent
word store, but it can break a whole-register comparison with a color assembled as
an unsigned value. Keep all cached and previous values in the same numerical form.
Extract individual components with logical shifts and an eight-bit mask. Store one
pixel with \texttt{sw}; advance its address by four with full-width pointer arithmetic.
